\section{Background}
Emotions are complex psychological states that manifest simultaneously across multiple channels: the musculature of the face, the prosody of the voice, and the electrochemical activity of the brain. The ability to decode these signals computationally---a discipline termed affective computing---has applications in clinical psychology, human-computer interaction, educational technology, and mental health monitoring. For most of its history, affective computing has framed its central task as a classification problem: given a set of observable cues, infer the emotion label that best describes the observed state \cite{Picard1997AffectiveComputing}.

This framing contains an embedded assumption that is rarely made explicit: that the expressed emotion and the experienced emotion are the same thing. In controlled laboratory conditions, where participants are instructed to display a specific emotion, this assumption is defensible. In naturalistic settings, and in clinically significant populations, it frequently is not. Humans regularly suppress, mask, or modify their emotional displays for social and contextual reasons. A patient in a clinical consultation may present as calm while experiencing significant distress. An individual with alexithymia---a condition characterised by difficulty identifying and expressing emotions---may display a neutral face while neurophysiological signals indicate elevated negative valence. A person with depression may exhibit socially appropriate affect while EEG recordings reveal attenuated reward-processing signatures in the theta band.

These observations motivate a reorientation of the affective computing problem. Rather than asking ``what emotion is being expressed?'', a more clinically meaningful question is: ``do the expressed emotion and the experienced emotion agree?'' This thesis develops a computational framework to answer the second question. The framework is built on the observation that electroencephalography (EEG), an involuntary neurophysiological signal driven by the autonomic nervous system, provides a partial window into the experienced emotional state, while audio and video signals capture the expressed emotional state, which is substantially under voluntary control. When the two agree, the affective response is coherent. When they diverge, the divergence is itself a signal worthy of characterisation---a signal we term cross-modal affective incoherence, or, from the perspective of the emitting system, emotion suppression.

The experimental vehicle for this investigation is the EAV dataset \cite{Lee2024EAV}, a laboratory corpus of 42 subjects with synchronised 30-channel EEG, audio, and video recordings across five emotion categories: Neutral, Anger, Happiness, Sadness, and Calmness. To our knowledge, this thesis presents the first characterisation of cross-modal affective coherence patterns on the EAV benchmark.

\section{Motivation}
The motivation for this research emerges from two converging observations: one architectural and one clinical.

The architectural observation is that the modalities available in systems like EAV are not epistemically equivalent. Audio and video are voluntary, external signals. A speaker chooses, at some level, how loud to speak, how much to smile, and what emotional register to adopt. EEG is an involuntary, internal signal. The power spectral density of the alpha, beta and theta bands reflects neurophysiological activity that is largely outside voluntary control, and that changes in response to affective states on timescales of tens to hundreds of milliseconds---far faster than the social-masking mechanisms that regulate facial expression and vocal prosody. Prior multimodal emotion recognition research has treated all three modalities as three noisy views of a single ground-truth emotional label, and has fused them accordingly. This thesis argues that a more precise model treats audio-video as the expressed label and EEG as the experienced label, and that the relationship between these two labels is the quantity of interest.

The clinical observation is that discordance between expressed and experienced emotion is not merely a theoretical curiosity. The emotion suppression literature documents the short- and long-term costs of suppression on psychological and physiological wellbeing \cite{Gross1998EmotionRegulation,Mauss2009Measures}. Disorders characterised by impaired emotional expressivity---alexithymia, depression, autism spectrum conditions, and post-traumatic stress disorder---are all associated with systematic discordance between self-reported emotional experience and observable affective behaviour. A computational tool that can detect and characterise this discordance from passively collected multimodal data has direct clinical relevance, provided the methodology is validated on population-scale data before clinical extrapolation is attempted.

The present thesis is explicitly scoped to the first step: developing the computational methodology and validating it on a healthy laboratory population. Clinical validation is framed as Phase 2 future work. This scoping decision, while limiting the direct clinical impact of the current submission, is the methodologically honest position given the constraints of a B.Sc. thesis and the laboratory nature of the EAV corpus.

\section{Problem Statement}
Current multimodal emotion recognition systems face three inter-related problems when applied to the goal of detecting cross-modal affective incoherence.

First, there is the modality-alignment problem. Multimodal fusion methods---whether early fusion through feature concatenation, late fusion through decision averaging, or attention-based fusion---all assume that the modalities should agree on a single label. When they are optimised to produce a consensus prediction, they suppress precisely the disagreement that the coherence-detection task requires to be detected. A system trained purely to maximise emotion classification accuracy will learn to resolve cross-modal conflicts, not to characterise them.

Second, there is the signal-quality problem in EEG. Electroencephalography in the context of natural emotional stimuli is a low signal-to-noise measurement. EEG is contaminated by electromyographic artifacts from facial muscles, motion artifacts from head movement, and between-session and between-subject non-stationarity. The effective per-subject training data in EAV is 280 trials---a regime in which the EEG encoder must learn to extract emotion-discriminative features from a 30-channel, 500-sample representation of five-second clips, without any cross-subject pretraining. This data scarcity means that EEG classifier confidence is a meaningful signal only when it is notably high, not in marginal cases. The suppression-event definition in this thesis therefore requires an EEG top-1 confidence threshold of 0.5, above which the EEG prediction is treated as a reliable indicator of internal state.

Third, there is the corpus-realism problem. The EAV corpus is a cue-based, laboratory-elicited dataset. Participants were given emotion cues and asked to respond; the emotions recorded are therefore elicited rather than spontaneous. This is not a defect of the corpus---it is the design choice that enables reliable ground-truth labelling---but it means that the suppression patterns observed in EAV are specifically those arising within the elicitation context, not the full range of suppression patterns observable in naturalistic discourse. The external validity of our findings is bounded by this corpus property.

These three problems collectively define the challenge this thesis addresses: building a trimodal classification system that is strong enough to extract reliable signal from each modality, and then using the resulting modality-level predictions to characterise the structure of cross-modal disagreement at population scale.

\section{Objectives and Research Questions}
The thesis addresses five formal research questions, each corresponding to a specific empirical contribution in the results section.

\begin{enumerate}
    \item \textbf{RQ1---Per-modality baseline.} Can each of the audio, video and EEG modalities be classified into the five EAV emotion categories at a rate that meaningfully exceeds the EAV paper's published per-modality baselines (audio: 36.7\%, vision: 52.8\%, EEG: 36.7\%)? A positive result is necessary to establish that each modality carries genuine discriminative signal, without which fusion and coherence analysis are not interpretable.
    \item \textbf{RQ2---Fusion gain.} Does multimodal fusion improve classification accuracy over the best single modality, and by what margin? The fusion gain quantifies how much complementary information the three modalities share and justifies the trimodal architecture over a strong single-modality baseline.
    \item \textbf{RQ3---Architectural contribution.} Does learned trimodal cross-modal attention fusion outperform a naive softmax-averaging late-fusion baseline at full 42-subject scale? This question tests whether the architectural investment in an attention-based fusion module is warranted, or whether a much simpler combination suffices on this dataset.
    \item \textbf{RQ4---Coherence as signal.} When the modalities disagree---specifically, when audio and video consensus diverges from EEG prediction---what patterns of disagreement emerge, and are those patterns consistent across the subject population? This is the thesis's primary research question, and its answer constitutes the headline empirical contribution.
    \item \textbf{RQ5---Robustness.} Does the system continue to produce meaningful predictions when the EEG modality is unavailable at inference time? This question is clinically necessary: any eventual deployment scenario will involve individuals who do not have an EEG headset, and the system must degrade gracefully rather than producing uninformative output.
\end{enumerate}

\section{Methodology in Brief}
The methodology proceeds in three stages, each building on the output of the previous.

\textbf{Stage 1: Per-modality training.} The Audio Spectrogram Transformer (AST), pre-trained on AudioSet, is fine-tuned on each subject's audio clips for 5 frozen epochs followed by 5 fine-tuning epochs. A Vision Transformer (ViT), pre-trained on a facial-emotion dataset, is fine-tuned on each subject's video frames for 3 frozen epochs followed by 2 fine-tuning epochs. EEGNet is trained from scratch on each subject's EEG segments for 350 epochs. All three models are trained within-subject using the EAV repository's stratified split (70\% train, 30\% test). Per-subject, per-modality test logits are saved as compressed NumPy archives for downstream consumption by Stages 2 and 3.

\textbf{Stage 2: Multimodal fusion.} The saved per-modality test logits are used in two ways. First, a naive late-fusion baseline is computed by mean-averaging the three per-modality softmax outputs and taking the argmax. Second, a trimodal cross-modal attention fusion module (\texttt{TrimodalAttentionFusion}) is trained on the per-subject, per-modality pre-classifier feature vectors. The fusion module projects each modality's features to a common 256-dimensional embedding, adds learned modality-type embeddings, passes the three-token sequence through a two-layer transformer encoder, mean-pools the output, and applies an MLP classification head.

\textbf{Stage 3: Coherence analysis.} The per-modality test predictions, not the fused predictions, are used to construct the suppression matrix. For each test trial, a suppression event is declared if audio and vision independently predict the same emotion label (external consensus), EEG predicts a different label (internal divergence), and EEG's top-1 confidence exceeds 0.5. Qualifying events are accumulated into a $5\times5$ contingency table indexed by external consensus emotion and EEG prediction, producing a population-level characterisation of cross-modal affective incoherence.

\section{Scope and Challenges}
\textbf{Scope.} This thesis is scoped to within-subject, clip-level emotion classification and coherence analysis on the EAV corpus. All 42 subjects are included. The five EAV emotion categories---Neutral, Anger, Happiness, Sadness, and Calmness---are the prediction targets. The inference granularity is one five-second clip per prediction. Cross-subject generalisation (LOSO evaluation) is included as a planned but partially executed component and is noted accordingly in the results. The demonstration component is a file-based video overlay using audio and video only, with EEG zeroed, consistent with any realistic deployment scenario.

\textbf{Challenges and resolutions.} The development of this thesis encountered thirteen distinct technical and methodological challenges, documented in detail in Chapter 7. The most significant of these were: (a) a forward-hook implementation defect in the EAV repository's EEGNet code that caused systematic training failure, resolved by extracting the hook into a named function with an implicit \texttt{None} return; (b) a train/eval mode cycling bug in the per-modality trainer, which suppressed BatchNorm learning from epoch 2 onwards, resolved by relocating the \texttt{model.train()} call inside the epoch loop; (c) a RAM-exhaustion failure in the vision preprocessing pipeline under eager frame loading, resolved by batched processing; and (d) the mid-project reformulation from accuracy-maximisation to coherence-characterisation, which required redefining the primary research question without discarding any of the existing code or experimental artefacts.
